\subsection{Marked Kirby--Cerf invariance and manifold descent}
\label{subsec:kirby-invariance}

The identity-character evaluator is already invariant under Kirby--Cerf
moves by \cref{rem:identity-character-specialization}.  We prove that the
same local ribbon identities remain valid after the character coupons are
inserted.  The only new inputs are the cocycle-coordinate transport of
\cref{prop:marking-coordinate-transport} and the naturality, monoidality,
and duality identities for \(\eta\).

\begin{lemma}[Character compatibility with elementary moves]
\label{lem:character-compatibility-elementary-moves}
Let
\[
  m\colon(\mathcal H,\beta)\longrightarrow(\mathcal H',\beta')
\]
be a marked elementary Kirby--Cerf move between one-ended Kirby data.  Choose
a cocycle \(\omega\) representing \(\xi_{\mathcal H,\beta}\), and choose the
target representative \(\omega'\) prescribed by
\cref{prop:marking-coordinate-transport}.  Every local ribbon identity used
to prove the identity-character invariance of \(m\) remains valid after the
source and target graphs are decorated by \(\omega\) and \(\omega'\),
respectively.
\end{lemma}

\begin{proof}
The nine directed move types available in a one-ended movie divide into
three kinds of character check.  We list every type explicitly.

\begin{itemize}
\item \textbf{Moves that do not change the relevant \(2\)-handle
coordinates.}
These are the following four move types.
\begin{itemize}
  \item \textbf{Level isotopy.}
  Every \(2\)-handle coordinate is unchanged.  Each character coupon moves
  through the same local morphisms as its ribbon, so naturality of \(\eta\)
  gives the identical target coupon.

  \item \textbf{A \(1/1\) slide.}
  Every \(2\)-handle coordinate is unchanged.  The untwisted
  dotted-projector identity accepts the already decorated framed link as its
  exterior input, so all character coupons pass through that identity
  unchanged.

  \item \textbf{A \(2/1\) passage through a dotted handle.}
  Every \(2\)-handle coordinate is unchanged.  The moving coupon is
  transported through the same local morphisms as its ribbon, and naturality
  of \(\eta\) gives the target coupon.

  \item \textbf{A \(3/3\) slide.}
  This changes only the basis of \(C_3\).  The dotted--framed graph and all
  \(C^2\)-coordinates are unchanged, hence so are the character coupons.
\end{itemize}

\item \textbf{\(2/2\) slides.}
Let \(x,y\) be the moving and stationary simple labels and let \(a,b\) be
their source cocycle coordinates.  The coordinate transport from
\cref{prop:marking-coordinate-transport} is
\[
  (a,b)\longmapsto(a+b,b)
  \quad\text{or}\quad
  (a,b)\longmapsto(a-b,b),
\]
according to the two orientations.  The two oriented local models are checked
separately:
\begin{itemize}
  \item \textbf{The positive \(2/2\) slide.}
  Its channel condition is \(z\subset x\otimes y^*\).  Monoidality and
  duality give
  \[
    \eta_x^{a+b}\eta_{z^*}^{b}
    =\eta_x^{a+b}(\eta_x\eta_y^{-1})^{-b}
    =\eta_x^a\eta_y^b,
  \]
  so its character factor is preserved.

  \item \textbf{The negative \(2/2\) slide.}
  Its channels satisfy \(z\subset x\otimes y\), and
\[
  \eta_x^{a-b}\eta_z^b
  =\eta_x^{a-b}(\eta_x\eta_y)^b
  =\eta_x^a\eta_y^b.
\]
  Thus its character factor is preserved as well.
\end{itemize}
Thus every fusion-channel term has exactly the same character factor before
and after attachment.  Naturality transports that factor along the band, and
the trace-dual detachment used in the untwisted proof preserves it.

\item \textbf{Adjacent-index births and deaths.}
These are the following four directed move types.
\begin{itemize}
  \item \textbf{A \(1\)--\(2\) birth.}
  The transported cohomology class has a representative
  \(\omega'=\omega\oplus0\), so the new coupon is
  \(\eta^0=\operatorname{id}\).

  \item \textbf{A \(1\)--\(2\) death.}
  Cohomological descent permits a representative whose disappearing
  \(2\)-handle coordinate is zero.  Its coupon is therefore the identity
  before the pair is removed.

  \item \textbf{A \(2\)--\(3\) birth.}
  The cocycle equation forces the new \(2\)-handle coordinate to be zero, so
  the new coupon is \(\eta^0=\operatorname{id}\).

  \item \textbf{A \(2\)--\(3\) death.}
  This is the reverse calculation: the cocycle equation forces the
  disappearing \(2\)-handle coordinate, and hence its character exponent,
  to be zero.
\end{itemize}
In all four cases the untwisted cancellation and normalization identities
therefore apply without alteration.
\end{itemize}
\end{proof}

\begin{theorem}[Elementary marked Kirby--Cerf invariance]
\label{thm:elementary-kirby-invariance}
Let
\[
  m\colon(\mathcal H,\beta)\longrightarrow(\mathcal H',\beta')
\]
be a marked elementary Kirby--Cerf move between one-ended Kirby data.  Then
\[
  \CY_{\cC,\eta}(\mathcal H,\beta)
  =\CY_{\cC,\eta}(\mathcal H',\beta').
\]
\end{theorem}

\begin{proof}
For the identity character this is the handle-decomposition invariance cited
in \cref{rem:identity-character-specialization}.  The proof is a finite set
of local ribbon identities, with the standard normalization changes for
adjacent-index cancellations.  By
\cref{lem:character-compatibility-elementary-moves}, every one of those local
identities remains valid for the transported character decorations.  By
\cref{prop:representative-independence}, the answer is independent of the
chosen source and target cocycle representatives.
\end{proof}

\begin{corollary}[Finite marked-movie invariance]
\label{cor:finite-marked-movie-invariance}
If a finite marked Kirby--Cerf movie connects
\((\mathcal H,\beta)\) to \((\mathcal K,\gamma)\), then
\[
  \CY_{\cC,\eta}(\mathcal H,\beta)
  =\CY_{\cC,\eta}(\mathcal K,\gamma).
\]
\end{corollary}

\begin{proof}
Apply \cref{thm:elementary-kirby-invariance} at each step and induct on the
length of the movie.
\end{proof}

\subsubsection{Descent to marked 4-manifolds}
\label{subsubsec:marked-manifold-descent}

\begin{proposition}[Well-definedness on a marked \(4\)-manifold]
\label{prop:marked-manifold-evaluation-well-defined}
Let
\[
  \rho\colon|(\mathcal H,\beta)|\longrightarrow(X,\alpha),
  \qquad
  \sigma\colon|(\mathcal K,\gamma)|\longrightarrow(X,\alpha)
\]
be marked one-ended Kirby presentations of the same marked \(4\)-manifold.
Then
\[
  \CY_{\cC,\eta}(\mathcal H,\beta)
  =\CY_{\cC,\eta}(\mathcal K,\gamma).
\]
\end{proposition}

\begin{proof}
Marked Kirby--Cerf completeness, \cref{prop:marked-kirby-cerf}, supplies a
finite marked movie between the two presentations.  Apply
\cref{cor:finite-marked-movie-invariance}.
\end{proof}

\begin{definition}[Manifold-level evaluator]
\label{def:manifold-level-evaluator}
For a marked \(4\)-manifold \((X,\alpha)\), choose a marked one-ended Kirby
presentation
\[
  \rho\colon|(\mathcal H,\beta)|\longrightarrow(X,\alpha)
\]
and define
\[
  \CY_{\cC,\eta}(X,\alpha)
  :=\CY_{\cC,\eta}(\mathcal H,\beta).
\]
This is independent of the chosen presentation by the preceding proposition.
\end{definition}

\begin{theorem}[Marked diffeomorphism invariance]
\label{prop:diffeomorphism-invariance}
Let
\[
  f\colon(X,\alpha)\longrightarrow(Y,\beta)
\]
be a marked diffeomorphism.  Then
\[
  \CY_{\cC,\eta}(X,\alpha)
  =\CY_{\cC,\eta}(Y,\beta).
\]
\end{theorem}

\begin{proof}
Choose marked one-ended Kirby presentations of the source and target.
Marked Kirby--Cerf completeness gives a finite marked movie realizing
\(f\), and \cref{cor:finite-marked-movie-invariance} identifies their
values.
\end{proof}
